\section{The Pelagian War, the Burned Monasteries, and the End}

\subsection{Last battle}

Jerome's final controversy joined him, at last, to the mainstream of Latin
theology. Pelagius, the British ascetic teacher whose followers held that
human nature can fulfill God's commands without inner grace and that a man
may be without sin if he will, had come east after the sack of Rome and
settled in Palestine, where John of Jerusalem received him kindly. Jerome
engaged in his own key. His letter to Ctesiphon (415) traced the new doctrine
to the philosophers' dream of passionlessness---quoting with relish ``one of
our own writers,'' Tertullian: ``the philosophers are the patriarchs of the
heretics''---and pressed the Pauline texts: if even the apostle does the evil
he would not, ``what becomes of the assertion that a man may be without sin
if he will?'' (\work{Letter} 133.1--2). The \work{Dialogue Against the
Pelagians} (415), his last major work of controversy, argued the same case
with unusual restraint, in the form of a conversation between an orthodox
and a Pelagian spokesman. At the synod of Diospolis in December 415 Pelagius
satisfied his Eastern judges; Jerome called it a miserable synod; Africa and
Rome eventually condemned the doctrine, with Augustine now leading and the
old man of Bethlehem cheering from his cell.

\subsection{The attack of 416}

Then the war came to his door. In 416 armed men fell on the monasteries of
Bethlehem. The most sober near-contemporary account is Augustine's, published
within a year or two: ``Certain followers of Pelagius are said to have most
cruelly beaten and maltreated the servants and handmaidens of the Lord who
lived under the care of the holy presbyter Jerome, slain his deacon, and
burnt his monastic houses; while he himself, by God's mercy, narrowly escaped
the violent attacks of these impious assailants in the shelter of a
well-defended fortress'' (\work{On the Proceedings of Pelagius} 66).
Augustine---note the ``are said''---declined to name the perpetrators and
left their punishment to the local bishops. Pope Innocent I did intervene, in
letters preserved in Jerome's own collection: he wrote to Jerome that he had
``hastened to put forth the authority of the apostolic see to repress the
plague in all its manifestations,'' but that ``as your letters name no
individuals and bring no specific charges, there is no one at present against
whom we can proceed''---if Jerome would lay a clear accusation, judges would
be appointed (\work{Letter} 136); a companion letter rebuked John of
Jerusalem for the outrage suffered in his diocese. Jerome's own surviving
notice is a single grim sentence to a correspondent the next year: ``As for
our house, so far as fleshly wealth is concerned, it has been completely
destroyed by the onslaughts of the heretics; but by the mercy of Christ it is
still filled with spiritual riches. To live on bread is better than to lose
the faith'' (\work{Letter} 139).

\witnesslimit{Responsibility for the attack is attested only as attribution:
Augustine's ``certain followers of Pelagius are said,'' Innocent's unnamed
``plague,'' Jerome's ``the heretics.'' No perpetrator was ever named in a
surviving source, no proceeding is recorded, and Pelagius himself was never
charged with ordering it. The event---beatings, a deacon killed, buildings
burned, Jerome besieged in a tower---is multiply attested; the
responsibility remains an accusation, and this study reports it as such.}

\subsection{Deaths}

The last years were a roll of losses. Paula had died in 404; Marcella,
beaten by the Goths in the sack of Rome---she ``showed no fear'' and pleaded
for her companion Principia---died of the aftermath in 410/411, mourned in
the letter that contains Jerome's epitaph on the city itself: ``The City
which had taken the whole world was itself taken'' (\work{Letter} 127.12--14).
Eustochium, shaken by the destruction of 416, died in 418/419. Jerome
survived them all, still dictating---his commentary on Jeremiah stopped
partway---and died at Bethlehem on 30 September of 419 or 420: the year 420
rests on the chronicle notice of Prosper of Aquitaine, the day on the
Church's immemorial commemoration, and official biographies (so Benedict
XVI's catechesis) give 419/420 without deciding. He was buried near Paula and
Eustochium, beside the cave of the Nativity. What later ages did with his
remains, his image, and his story belongs to the tradition history examined
next.
